\documentclass[12pt]{article}

\usepackage[T1]{fontenc}
\usepackage[utf8]{inputenc}
\usepackage[margin=1in]{geometry}
\usepackage{amsmath,amssymb,booktabs,array,xcolor,makecell}
\usepackage{newpxtext,newpxmath}
\usepackage{enumitem}
\usepackage{fancyhdr}
\usepackage{parskip}

\pagestyle{fancy}
\fancyhf{}
\lhead{\textit{ECON 1001 \textperiodcentered{} Spring 2026}}
\rhead{\textit{Final Exam Solutions}}
\cfoot{\thepage}
\renewcommand{\headrulewidth}{0.4pt}

\setlength{\parskip}{0.5em}
\setlength{\parindent}{0pt}

\newlist{subqs}{enumerate}{1}
\setlist[subqs]{
  label=\textbf{\alph*)},
  leftmargin=1.8em,
  itemsep=0.6em,
  topsep=0.3em
}

\newlist{solbullets}{itemize}{1}
\setlist[solbullets]{
  label=\textbullet,
  leftmargin=1.4em,
  itemsep=1pt,
  topsep=3pt,
  parsep=0pt
}

\begin{document}

% ── TITLE ─────────────────────────────────────────────────────────────────────
\thispagestyle{empty}
\begin{center}
  {\large\bfseries ECON 1001 \textperiodcentered{} Principles of Economics \textperiodcentered{} Spring 2026}\\[6pt]
  {\LARGE\bfseries Final Exam --- Solutions}\\[6pt]
  {\normalsize\itshape Your exam version (A--E) is printed on the cover page of your exam.}
\end{center}

\vspace{0.8em}
\noindent\rule{\linewidth}{0.6pt}
\vspace{0.5em}

% ── PART I: MC ANSWER KEY ─────────────────────────────────────────────────────
\section*{Part I\quad Multiple Choice Answer Key (Problems 1--27)}

Find your version column and look up each question by its number on your exam (1--27). The correct answer for that question on your version is shown in the corresponding cell.

\vspace{0.6em}
{\small
\begin{center}
\renewcommand{\arraystretch}{1.08}
\setlength{\tabcolsep}{5.5pt}
\begin{tabular}{r ccccc @{\hspace{2em}} r ccccc}
  \toprule
  \textbf{\#} & \textbf{Ver.\ A} & \textbf{Ver.\ B} & \textbf{Ver.\ C} & \textbf{Ver.\ D} & \textbf{Ver.\ E} &
  \textbf{\#} & \textbf{Ver.\ A} & \textbf{Ver.\ B} & \textbf{Ver.\ C} & \textbf{Ver.\ D} & \textbf{Ver.\ E} \\
  \midrule
  \textbf{1}  & D & B & C & C & C &  \textbf{15} & B & B & B & C & D \\
  \textbf{2}  & D & D & D & D & C &  \textbf{16} & D & B & B & C & D \\
  \textbf{3}  & D & C & A & B & A &  \textbf{17} & B & C & A & C & B \\
  \textbf{4}  & D & D & C & C & D &  \textbf{18} & A & A & B & C & D \\
  \textbf{5}  & D & C & B & C & D &  \textbf{19} & D & A & A & A & B \\
  \textbf{6}  & A & C & C & C & D &  \textbf{20} & B & B & C & C & D \\
  \textbf{7}  & A & D & C & A & B &  \textbf{21} & D & B & B & A & B \\
  \textbf{8}  & A & A & D & A & D &  \textbf{22} & D & B & C & D & C \\
  \textbf{9}  & B & B & A & D & C &  \textbf{23} & C & A & A & B & D \\
  \textbf{10} & C & D & C & C & A &  \textbf{24} & D & C & A & D & C \\
  \textbf{11} & C & B & D & D & B &  \textbf{25} & B & D & B & A & D \\
  \textbf{12} & A & D & D & C & C &  \textbf{26} & D & A & B & A & A \\
  \textbf{13} & D & C & A & A & B &  \textbf{27} & C & B & D & D & C \\
  \textbf{14} & C & B & B & A & A &              &   &   &   &   &   \\
  \bottomrule
\end{tabular}
\end{center}
}

\vspace{1em}
\noindent\rule{\linewidth}{0.6pt}
\vspace{0.5em}

% ── PART II: FR SOLUTIONS ─────────────────────────────────────────────────────
\section*{Part II\quad Free Response Solutions (Problem 28)}

% ── VERSION A ─────────────────────────────────────────────────────────────────
\subsection*{Version A}

\textbf{Parameters:} Aggregate Demand: $P = 10 - 0.5Q$ \textbar{} $MC = 4$

\begin{subqs}

\item \textbf{Single Price \& Quantity}
  \begin{solbullets}
    \item Set $MR = MC$ to find $Q^*$.
    \item \textbf{$Q^* = 6$}
    \item Plug $Q^*$ into Aggregate Demand to find $P^*$.
    \item \textbf{$P^* = \$7$}
  \end{solbullets}

\item \textbf{Deadweight Loss (DWL)}
  \begin{solbullets}
    \item Efficient Quantity ($Q_{eff}$) is where $P = MC \rightarrow Q = 12$.
    \item $DWL = \frac{1}{2} \times (P^* - MC) \times (Q_{eff} - Q^*)$
    \item $DWL = 0.5 \times (7 - 4) \times (12 - 6)$
    \item \textbf{$DWL = 9$}
  \end{solbullets}

\item \textbf{Total Profit at Single Price}
  \begin{solbullets}
    \item Profit $= (P^* - MC) \times Q^*$
    \item Profit $= (7 - 4) \times 6$
    \item \textbf{Profit $= \$18$}
  \end{solbullets}

\item \textbf{Price if firm CANNOT distinguish groups}
  \begin{solbullets}
    \item If the firm cannot distinguish, they face the aggregate demand curve.
    \item They will charge the same single price found in part (a).
    \item \textbf{Price $= \$7$}
  \end{solbullets}

\item \textbf{Price Discrimination (Distinguishing Groups)}
  \begin{solbullets}
    \item Act as a monopolist in each market: Set $MR = MC$ for Locals and Visitors.
    \item \textbf{Locals:} $P_L = 12 - Q_L \rightarrow$ Set $MR_L = MC \rightarrow Q_L = 4 \rightarrow$ \textbf{$P_L = \$8$}
    \item \textbf{Visitors:} $P_V = 8 - Q_V \rightarrow$ Set $MR_V = MC \rightarrow Q_V = 2 \rightarrow$ \textbf{$P_V = \$6$}
    \item \textbf{New Total Profit:} $(P_L - MC) \times Q_L + (P_V - MC) \times Q_V = (8 - 4) \times 4 + (6 - 4) \times 2$
    \item \textbf{Profit $= \$20$}
  \end{solbullets}

\item \textbf{Does the socially efficient price change?}
  \begin{solbullets}
    \item \textbf{No.} The socially efficient price is always where Marginal Benefit equals Marginal Cost ($P = MC$).
    \item Because the marginal cost remains $\$4$ for both groups, the socially efficient price remains $\$4$. Price discrimination changes the \textit{profit-maximizing} prices, but not the \textit{socially efficient} outcome.
  \end{solbullets}

\end{subqs}

\newpage

% ── VERSION B ─────────────────────────────────────────────────────────────────
\subsection*{Version B}

\textbf{Parameters:} Aggregate Demand: $P = 12 - 0.5Q$ \textbar{} $MC = 4$

\begin{subqs}

\item \textbf{Single Price \& Quantity}
  \begin{solbullets}
    \item Set $MR = MC$ to find $Q^*$.
    \item \textbf{$Q^* = 8$}
    \item Plug $Q^*$ into Aggregate Demand to find $P^*$.
    \item \textbf{$P^* = \$8$}
  \end{solbullets}

\item \textbf{Deadweight Loss (DWL)}
  \begin{solbullets}
    \item Efficient Quantity ($Q_{eff}$) is where $P = MC \rightarrow Q = 16$.
    \item $DWL = \frac{1}{2} \times (P^* - MC) \times (Q_{eff} - Q^*)$
    \item $DWL = 0.5 \times (8 - 4) \times (16 - 8)$
    \item \textbf{$DWL = 16$}
  \end{solbullets}

\item \textbf{Total Profit at Single Price}
  \begin{solbullets}
    \item Profit $= (P^* - MC) \times Q^*$
    \item Profit $= (8 - 4) \times 8$
    \item \textbf{Profit $= \$32$}
  \end{solbullets}

\item \textbf{Price if firm CANNOT distinguish groups}
  \begin{solbullets}
    \item If the firm cannot distinguish, they face the aggregate demand curve.
    \item They will charge the same single price found in part (a).
    \item \textbf{Price $= \$8$}
  \end{solbullets}

\item \textbf{Price Discrimination (Distinguishing Groups)}
  \begin{solbullets}
    \item Act as a monopolist in each market: Set $MR = MC$ for Locals and Visitors.
    \item \textbf{Locals:} $P_L = 14 - Q_L \rightarrow$ Set $MR_L = MC \rightarrow Q_L = 5 \rightarrow$ \textbf{$P_L = \$9$}
    \item \textbf{Visitors:} $P_V = 10 - Q_V \rightarrow$ Set $MR_V = MC \rightarrow Q_V = 3 \rightarrow$ \textbf{$P_V = \$7$}
    \item \textbf{New Total Profit:} $(P_L - MC) \times Q_L + (P_V - MC) \times Q_V = (9 - 4) \times 5 + (7 - 4) \times 3$
    \item \textbf{Profit $= \$34$}
  \end{solbullets}

\item \textbf{Does the socially efficient price change?}
  \begin{solbullets}
    \item \textbf{No.} The socially efficient price is always where Marginal Benefit equals Marginal Cost ($P = MC$).
    \item Because the marginal cost remains $\$4$ for both groups, the socially efficient price remains $\$4$. Price discrimination changes the \textit{profit-maximizing} prices, but not the \textit{socially efficient} outcome.
  \end{solbullets}

\end{subqs}

\newpage

% ── VERSION C ─────────────────────────────────────────────────────────────────
\subsection*{Version C}

\textbf{Parameters:} Aggregate Demand: $P = 16 - Q$ \textbar{} $MC = 4$

\begin{subqs}

\item \textbf{Single Price \& Quantity}
  \begin{solbullets}
    \item Set $MR = MC$ to find $Q^*$.
    \item \textbf{$Q^* = 6$}
    \item Plug $Q^*$ into Aggregate Demand to find $P^*$.
    \item \textbf{$P^* = \$10$}
  \end{solbullets}

\item \textbf{Deadweight Loss (DWL)}
  \begin{solbullets}
    \item Efficient Quantity ($Q_{eff}$) is where $P = MC \rightarrow Q = 12$.
    \item $DWL = \frac{1}{2} \times (P^* - MC) \times (Q_{eff} - Q^*)$
    \item $DWL = 0.5 \times (10 - 4) \times (12 - 6)$
    \item \textbf{$DWL = 18$}
  \end{solbullets}

\item \textbf{Total Profit at Single Price}
  \begin{solbullets}
    \item Profit $= (P^* - MC) \times Q^*$
    \item Profit $= (10 - 4) \times 6$
    \item \textbf{Profit $= \$36$}
  \end{solbullets}

\item \textbf{Price if firm CANNOT distinguish groups}
  \begin{solbullets}
    \item If the firm cannot distinguish, they face the aggregate demand curve.
    \item They will charge the same single price found in part (a).
    \item \textbf{Price $= \$10$}
  \end{solbullets}

\item \textbf{Price Discrimination (Distinguishing Groups)}
  \begin{solbullets}
    \item Act as a monopolist in each market: Set $MR = MC$ for Locals and Visitors.
    \item \textbf{Locals:} $P_L = 20 - 2Q_L \rightarrow$ Set $MR_L = MC \rightarrow Q_L = 4 \rightarrow$ \textbf{$P_L = \$12$}
    \item \textbf{Visitors:} $P_V = 12 - 2Q_V \rightarrow$ Set $MR_V = MC \rightarrow Q_V = 2 \rightarrow$ \textbf{$P_V = \$8$}
    \item \textbf{New Total Profit:} $(P_L - MC) \times Q_L + (P_V - MC) \times Q_V = (12 - 4) \times 4 + (8 - 4) \times 2$
    \item \textbf{Profit $= \$40$}
  \end{solbullets}

\item \textbf{Does the socially efficient price change?}
  \begin{solbullets}
    \item \textbf{No.} The socially efficient price is always where Marginal Benefit equals Marginal Cost ($P = MC$).
    \item Because the marginal cost remains $\$4$ for both groups, the socially efficient price remains $\$4$. Price discrimination changes the \textit{profit-maximizing} prices, but not the \textit{socially efficient} outcome.
  \end{solbullets}

\end{subqs}

\newpage

% ── VERSION D ─────────────────────────────────────────────────────────────────
\subsection*{Version D}

\textbf{Parameters:} Aggregate Demand: $P = 10 - 0.25Q$ \textbar{} $MC = 4$

\begin{subqs}

\item \textbf{Single Price \& Quantity}
  \begin{solbullets}
    \item Set $MR = MC$ to find $Q^*$.
    \item \textbf{$Q^* = 12$}
    \item Plug $Q^*$ into Aggregate Demand to find $P^*$.
    \item \textbf{$P^* = \$7$}
  \end{solbullets}

\item \textbf{Deadweight Loss (DWL)}
  \begin{solbullets}
    \item Efficient Quantity ($Q_{eff}$) is where $P = MC \rightarrow Q = 24$.
    \item $DWL = \frac{1}{2} \times (P^* - MC) \times (Q_{eff} - Q^*)$
    \item $DWL = 0.5 \times (7 - 4) \times (24 - 12)$
    \item \textbf{$DWL = 18$}
  \end{solbullets}

\item \textbf{Total Profit at Single Price}
  \begin{solbullets}
    \item Profit $= (P^* - MC) \times Q^*$
    \item Profit $= (7 - 4) \times 12$
    \item \textbf{Profit $= \$36$}
  \end{solbullets}

\item \textbf{Price if firm CANNOT distinguish groups}
  \begin{solbullets}
    \item If the firm cannot distinguish, they face the aggregate demand curve.
    \item They will charge the same single price found in part (a).
    \item \textbf{Price $= \$7$}
  \end{solbullets}

\item \textbf{Price Discrimination (Distinguishing Groups)}
  \begin{solbullets}
    \item Act as a monopolist in each market: Set $MR = MC$ for Locals and Visitors.
    \item \textbf{Locals:} $P_L = 12 - 0.5Q_L \rightarrow$ Set $MR_L = MC \rightarrow Q_L = 8 \rightarrow$ \textbf{$P_L = \$8$}
    \item \textbf{Visitors:} $P_V = 8 - 0.5Q_V \rightarrow$ Set $MR_V = MC \rightarrow Q_V = 4 \rightarrow$ \textbf{$P_V = \$6$}
    \item \textbf{New Total Profit:} $(P_L - MC) \times Q_L + (P_V - MC) \times Q_V = (8 - 4) \times 8 + (6 - 4) \times 4$
    \item \textbf{Profit $= \$40$}
  \end{solbullets}

\item \textbf{Does the socially efficient price change?}
  \begin{solbullets}
    \item \textbf{No.} The socially efficient price is always where Marginal Benefit equals Marginal Cost ($P = MC$).
    \item Because the marginal cost remains $\$4$ for both groups, the socially efficient price remains $\$4$. Price discrimination changes the \textit{profit-maximizing} prices, but not the \textit{socially efficient} outcome.
  \end{solbullets}

\end{subqs}

\newpage

% ── VERSION E ─────────────────────────────────────────────────────────────────
\subsection*{Version E}

\textbf{Parameters:} Aggregate Demand: $P = 20 - 0.5Q$ \textbar{} $MC = 8$

\begin{subqs}

\item \textbf{Single Price \& Quantity}
  \begin{solbullets}
    \item Set $MR = MC$ to find $Q^*$.
    \item \textbf{$Q^* = 12$}
    \item Plug $Q^*$ into Aggregate Demand to find $P^*$.
    \item \textbf{$P^* = \$14$}
  \end{solbullets}

\item \textbf{Deadweight Loss (DWL)}
  \begin{solbullets}
    \item Efficient Quantity ($Q_{eff}$) is where $P = MC \rightarrow Q = 24$.
    \item $DWL = \frac{1}{2} \times (P^* - MC) \times (Q_{eff} - Q^*)$
    \item $DWL = 0.5 \times (14 - 8) \times (24 - 12)$
    \item \textbf{$DWL = 36$}
  \end{solbullets}

\item \textbf{Total Profit at Single Price}
  \begin{solbullets}
    \item Profit $= (P^* - MC) \times Q^*$
    \item Profit $= (14 - 8) \times 12$
    \item \textbf{Profit $= \$72$}
  \end{solbullets}

\item \textbf{Price if firm CANNOT distinguish groups}
  \begin{solbullets}
    \item If the firm cannot distinguish, they face the aggregate demand curve.
    \item They will charge the same single price found in part (a).
    \item \textbf{Price $= \$14$}
  \end{solbullets}

\item \textbf{Price Discrimination (Distinguishing Groups)}
  \begin{solbullets}
    \item Act as a monopolist in each market: Set $MR = MC$ for Locals and Visitors.
    \item \textbf{Locals:} $P_L = 24 - Q_L \rightarrow$ Set $MR_L = MC \rightarrow Q_L = 8 \rightarrow$ \textbf{$P_L = \$16$}
    \item \textbf{Visitors:} $P_V = 16 - Q_V \rightarrow$ Set $MR_V = MC \rightarrow Q_V = 4 \rightarrow$ \textbf{$P_V = \$12$}
    \item \textbf{New Total Profit:} $(P_L - MC) \times Q_L + (P_V - MC) \times Q_V = (16 - 8) \times 8 + (12 - 8) \times 4$
    \item \textbf{Profit $= \$80$}
  \end{solbullets}

\item \textbf{Does the socially efficient price change?}
  \begin{solbullets}
    \item \textbf{No.} The socially efficient price is always where Marginal Benefit equals Marginal Cost ($P = MC$).
    \item Because the marginal cost remains $\$8$ for both groups, the socially efficient price remains $\$8$. Price discrimination changes the \textit{profit-maximizing} prices, but not the \textit{socially efficient} outcome.
  \end{solbullets}

\end{subqs}

\end{document}
